% =====================================================================
%  THÈME 2 — Conversions en quantité de matière
%  Niveau FACILE — Corrigé
% =====================================================================
\documentclass[11pt,a4paper]{article}
\usepackage[utf8]{inputenc}
\usepackage[T1]{fontenc}
\usepackage{lmodern}
\usepackage[french]{babel}
\usepackage{amsmath,amssymb,mathtools}
\usepackage{icomma}
\usepackage[version=4]{mhchem}
\usepackage{siunitx}
\sisetup{output-decimal-marker={,},per-mode=symbol}
\DeclareSIUnit\Molar{\mole\per\litre}
\usepackage{xcolor}
\usepackage{tikz}
\usepackage[most]{tcolorbox}
\usepackage{enumitem}
\usepackage{array}
\usepackage{fancyhdr}
\usepackage[a4paper,margin=2.1cm,headheight=26pt,headsep=8pt]{geometry}

\definecolor{primary}{RGB}{142,93,168}
\definecolor{primarydark}{RGB}{82,45,108}
\definecolor{corrcol}{RGB}{176,110,20}
\newcommand{\gmol}{\si{\gram\per\mole}}
\newcommand{\Vm}{V_{\mathrm m}}

\pagestyle{fancy}\fancyhf{}
\fancyhead[L]{\small\textcolor{primarydark}{\textbf{Fiche 5} — Thème 2 : Conversions}}
\fancyhead[R]{\small\textcolor{corrcol}{\textsc{Corrigé — Facile}}}
\fancyfoot[C]{\small\thepage}
\renewcommand{\headrulewidth}{0.8pt}
\renewcommand{\headrule}{\hbox to\headwidth{\color{primary}\leaders\hrule height \headrulewidth\hfill}}

\newtcolorbox[auto counter]{cor}[1][]{enhanced,breakable,boxrule=1pt,arc=3pt,
  left=3mm,right=3mm,top=2.4mm,bottom=2.4mm,colback=corrcol!6,colframe=corrcol,
  fonttitle=\bfseries,coltitle=white,
  title={Corrigé~\thetcbcounter\ \ifx\relax#1\relax\relax\else\textmd{--- \itshape #1}\fi}}

\newcommand{\rep}[1]{\par\smallskip\noindent\textbf{\color{corrcol}$\Rightarrow$ #1}}

\setlength{\parindent}{0pt}
\setlength{\parskip}{3pt}

\begin{document}

\begin{center}
{\LARGE\bfseries\color{primarydark}Thème 2 — Conversions}\\[1pt]
{\large\color{corrcol}\textbf{Corrigé — Niveau Facile}}\\
{\color{primary}\rule{0.6\linewidth}{1pt}}
\end{center}

\begin{cor}[calculer une masse molaire]
\textbf{a)} $M(\ce{H2SO4}) = 2\times1{,}0 + 1\times32{,}0 + 4\times16{,}0 = 2+32+64 = \boxed{\SI{98}{\gram\per\mole}}$.\\
\textbf{b)} $M(\ce{CaCO3}) = 40{,}0 + 12{,}0 + 3\times16{,}0 = 40+12+48 = \boxed{\SI{100}{\gram\per\mole}}$.
\end{cor}

\begin{cor}[masse $\leftrightarrow$ mole]
\textbf{a)} $n=\dfrac{m}{M}=\dfrac{20{,}0}{40{,}0}=\boxed{\SI{0,50}{\mole}}$.\\
\textbf{b)} On inverse la relation : $m=n\times M = 0{,}25\times180 = \boxed{\SI{45}{\gram}}$.
\rep{Réflexe : « combien de moles ? » $\to \div M$ ; « quelle masse ? » $\to \times M$.}
\end{cor}

\begin{cor}[volume d'un gaz aux TPN]
\textbf{a)} $V = n\times\Vm = 0{,}25\times22{,}4 = \boxed{\SI{5,6}{\litre}}$.\\
\textbf{b)} $n = \dfrac{V}{\Vm} = \dfrac{44{,}8}{22{,}4} = \boxed{\SI{2,0}{\mole}}$.
\rep{Aux TPN seulement : $\Vm=\SI{22,4}{\litre\per\mole}$.}
\end{cor}

\begin{cor}[quantité de matière en solution]
\textbf{a)} $\SI{250}{\milli\litre} = \SI{0,250}{\litre}$ (on divise par 1000).\\
\textbf{b)} $n = C\times V = 0{,}20\times0{,}250 = \boxed{\SI{0,050}{\mole}}$.
\rep{Piège classique : si on avait gardé $V=250$, on trouverait \SI{50}{\mole}, absurde. Toujours convertir mL $\to$ L.}
\end{cor}

\begin{cor}[compter les entités avec $N_A$]
\textbf{a)} $N = n\times N_A = 0{,}10\times\SI{6,02e23}{} = \boxed{\SI{6,02e22}{}}$ molécules.\\
\textbf{b)} $n = \dfrac{N}{N_A} = \dfrac{\SI{3,01e23}{}}{\SI{6,02e23}{}} = \boxed{\SI{0,50}{\mole}}$.
\end{cor}

\end{document}
